\Needspace{12\baselineskip}
\originalchapter{旋转与沿岸流动}
\originalsection{旋转坐标和Taylor--Proudman约束}
转速 $\boldsymbol\Omega$ 的坐标系中，动量含Coriolis项 $2\boldsymbol\Omega\times\mathbf u$；离心势通常并入有效重力和压力。局部平面设 $f=2\Omega\sin\varphi$，地转近似 $f\hat z\times\mathbf u=-\nabla_h p/\rho$，适用 $Ro\ll1$、演化慢、远离强摩擦边界的条件。

稳态、均匀密度、无黏、弱惯性且转速常量时，对地转平衡取旋度，结合不可压缩得 $\partial_z\mathbf u=0$，即Taylor--Proudman柱状约束。密度分层可以给热成风剪切，边界摩擦也破坏柱状性；不能将它概括为所有海洋流都与深度无关。

\originalsection{Ekman层}
地转外流附近的水平扰动 $u,v$ 满足 $-fv=\nu u_{zz}$、$fu=\nu v_{zz}$。用复速度 $W=u+\ii v$ 合并为 $\nu W_{zz}=\ii fW$，衰减根含实虚部分，速度随深度衰减并旋转，尺度 $\delta_E=\sqrt{2\nu/|f|}$。风应力边界或底部无滑移边界产生不同振幅和方向。

表面风应力的积分输运可从深度积分动量得到 $f\hat z\times\mathbf M=\boldsymbol\tau/\rho$（深部应力忽略），所以 $\mathbf M=-\hat z\times\boldsymbol\tau/(\rho f)$。输运散度由连续方程驱使垂直速度，沿岸边界会导致上升流或下沉流。符号须结合半球、岸线法向和风向检查，不能背一句“总向右”代替坐标定义。

\originalsection{浅水波和两层海洋}
深度 $H+\eta$、水平速度 $\mathbf u_h$ 的浅水方程把垂向结构平均，要求波长远大于水深、压力近静水。在线性、无背景流时
\[
u_t-fv=-g\eta_x,\quad v_t+fu=-g\eta_y,\quad \eta_t+H(u_x+v_y)=0.
\]
平面波给 $\omega^2=f^2+gHk^2$。忽略旋转后速度 $c=\sqrt{gH}$；旋转引入最低惯性频率和变形半径 $R_D=\sqrt{gH}/|f|$。

沿直岸的Kelvin波可取岸法向速度零，由横向地转平衡得到离岸指数衰减，沿岸相速仍为 $\sqrt{gH}$，只允许与半球相适配的一种传播方向。两层密度接近的海洋还有内部模，约化重力 $g'=g(\rho_2-\rho_1)/\rho_0$，长波相速 $c_i^2=g'H_1H_2/(H_1+H_2)$。沿岸风可以同时激发快的外模和慢的内模，需区分表面高度与内部界面位移。源课具体旋风、瞬变风和两层上升流的完整初边值解不在本版逐式复现范围。

\originalsection{旋转变换和小Rossby层次}
绝对速度 $\mathbf u_a=\mathbf u+\boldsymbol\Omega\times\mathbf r$。对向量时间导数再作旋转变换，恒定转速时加速度多 $2\boldsymbol\Omega\times\mathbf u+\boldsymbol\Omega\times(\boldsymbol\Omega\times\mathbf r)$；若转速变化另有Euler项。离心项为势梯度，可进入静水背景，Coriolis与速度垂直不直接做功，却改变方向和能量传输路径。

$Ro=U/(fL)\ll1$ 不自动保证稳态；若时间尺度 $T$ 与 $1/f$ 相近，局部加速度仍重要，产生惯性振荡。远离边界、慢变化的领先阶才是地转。近岸边界和Ekman层使摩擦与Coriolis同阶，形成不同近似区域。

\originalsection{热成风与柱状约束的破坏}
地转加静水 $f\hat z\times\mathbf u_h=-\nabla_hp/\rho_0,p_z=-\rho g$，对深度微分得
\[
f\hat z\times\partial_z\mathbf u_h=\frac g{\rho_0}\nabla_h\rho,
\quad u_z=\frac g{f\rho_0}\rho_y,
\quad v_z=-\frac g{f\rho_0}\rho_x.
\]
横向密度梯度允许竖向剪切，是热成风关系。密度均匀时才恢复Taylor--Proudman柱状性。地转速度由压力梯度决定，绝对深层参考流往往还需额外观测；密度场只能给剪切，并不无条件确定全部速度。

\originalsection{Ekman螺旋与输运积分}
深海向下距表面 $z_d>0$，相对地转速度的复扰动 $W$ 满足 $\nu W''=ifW$，$f>0$ 衰减根 $(1+i)/\delta_E$。表面应力边界选定复振幅，底层无滑移则选另一振幅。源课允许用恒定涡黏性近似，但现实湍流系数随深度变化，理想螺旋不是每次观测都完整出现。

将水平动量对深度积分，表面应力而深层剪切为零时有 $f\hat z\times\mathbf M=\boldsymbol\tau/\rho_0$。对 $f>0$，向北风应力 $\tau_y>0$ 给向东输运 $M_x=\tau_y/(\rho_0f)$。若海在西岸的东侧，输运离岸，连续式要求深水补上，产生上升流。这个坐标例子将半球与海岸方向都写明。

空间变化风应力产生 $w_E=\nabla_h\cdot\mathbf M$（符号按垂向坐标约定）；在常 $f$ 区等价于风应力旋度除 $\rho_0f$。若 $f$ 随纬度变化，散度还含其导数。底摩擦和沿岸压力梯度也会改变总水柱输运，不能把Ekman输运当成整个流量。

\originalsection{浅水势涡量、惯性重力与Rossby波}
无摩擦浅水质量守恒 $D(H+\eta)/Dt=-(H+\eta)\nabla_h\cdot\mathbf u$。对动量取旋度，绝对涡量满足 $D(\zeta+f)/Dt=-(\zeta+f)\nabla_h\cdot\mathbf u$，相除得 $D[(\zeta+f)/(H+\eta)]/Dt=0$。水柱被拉薄/拉厚必须相应调整旋转，这是势涡量守恒。

线性常 $f$ 方程的惯性重力波给 $\omega^2=f^2+c^2(k^2+l^2)$，$c=\sqrt{gH}$；还有零频地转分支。缓慢变化 $f=f_0+\beta y$、近地转流用 $u=-\psi_y,v=\psi_x$，领先阶势涡量异常 $q'=\nabla_h^2\psi-\psi/R_D^2$。绕静止基本态线性化给
\[
\partial_tq'+\beta\psi_x=0,\qquad
\omega=-\frac{\beta k}{k^2+l^2+R_D^{-2}}.
\]
这是长尺度Rossby波，传播方向与高频重力模不同。加入背景流时要同时修改背景势涡量梯度，不能只任意加一个平移速度而保留不一致的水面坡度。

\originalsection{岸边Kelvin波的方向和两层模}
取直岸 $y=0$、海在 $y>0$、北半球 $f>0$，岸法向速度 $v=0$。沿岸动量和连续式给相速 $c$，横向地转给 $\eta_y=-(f/c)\eta$，所以 $\eta\propto\e^{-y/R_D}\e^{ik(x-ct)}$。衰减选取 $c>0$，即在这个坐标中向 $+x$，海岸在传播右侧。改变海的一侧或半球需重做符号，而非背诵一个不带坐标的方位。

两层系统有自由面快外模和界面慢内模。密度差小且表面近刚盖的内部长波给 $c_i^2=g'H_1H_2/(H_1+H_2)$，变形半径 $R_i=c_i/|f|$。沿岸持续风使上层Ekman离岸，界面升高，内部Kelvin/重力信号沿岸和离岸传送调整；初始阶段、强迫持续阶段和停风后自由调整是不同初边值问题。

源课的旋风、瞬变大尺度风、两层与沿岸上升流讲义实际存在。本文给出共同方程、模、边界和输运机制；具体给定风场的解析积分与每幅原图逐式复现仍需额外篇幅，未伪称已有完整数值海洋模拟。
